% Finite spectral windows and the massless electroweak limit: short note, from lemmas 28, 29 and 30
\ifdefined\mobile
\documentclass[aps,prd,reprint,onecolumn,nofootinbib,amsmath,amssymb,floatfix]{revtex4-2}
\usepackage[paperwidth=270pt,paperheight=585pt,textwidth=246pt,textheight=545pt,top=18pt,left=12pt,headheight=0pt,headsep=0pt,footskip=14pt]{geometry}
\else
\documentclass[aps,prd,reprint,nofootinbib,amsmath,amssymb,floatfix]{revtex4-2}
\fi
\usepackage[T1]{fontenc}
\usepackage{lmodern}
\usepackage{amsthm}
\emergencystretch=1.5em
\ifdefined\mobile\makeatletter\def\frontmatter@abstractwidth{246\p@}\makeatother\fi
\usepackage[hidelinks]{hyperref}
\theoremstyle{plain}
\newtheorem{theorem}{Theorem}
\newtheorem{proposition}[theorem]{Proposition}
\newtheorem{corollary}[theorem]{Corollary}
\newcommand{\C}{\mathbb C}
\renewcommand{\H}{\mathbb H}
\newcommand{\R}{\mathbb R}
\newcommand{\CP}{\mathbb{CP}}
\newcommand{\hatD}{\widehat D}
\newcommand{\supp}[1]{Supplemental Material, part #1}

\begin{document}

\title{Finite spectral windows and the massless electroweak limit of seven-dimensional parents}
\date{Draft, 4 November 2026}

\begin{abstract}
The statement that a Kaluza-Klein dimension opens as the $W$ and $Z$ become massless needs a family, since along a general path to $M_W=M_Z=0$ no tower opens. We give families with the hierarchy $m\ll E_\varepsilon\ll\Delta\to0$: a moving low-energy window keeps a fixed finite number of states, and a fixed positive energy resolves more and more of the parent. On an interval this follows from exact formulas. On the disc bundle of $\Lambda^2_-T^*S^4$, a connected seven-dimensional commutative parent with boundary $\CP^3$, we prove for a chosen class of metrics that a local Dirac operator with a Callias potential has exactly eight light states, one copy of the $\mathbf 4$ of $Sp(2)$ at each end, with exponentially small mixing, and that real functions compress to the two-sheet finite maps generating $M_2(\H)\oplus M_4(\C)$. Along a family with logarithmically growing potential and expanding metric, the eight states stay separated at the moving scale, and the spectrum at fixed energy has seven-dimensional Weyl growth. The proofs rest on exact sums of squares given by $Sp(2)$-equivariance. The vector poles on the disc bundle, the gauge couplings, the connector of the real bimodule and any dynamics selecting the family remain open.
\end{abstract}

\maketitle

\section{Introduction}

In the towerless reading of almost-commutative geometry \cite{towerlesspaper}, a finite spectral triple is the low-energy window of a parent with an internal geometry, and the labels 9 and 11 are limits of an even interior at finite Fermi constant, where the physical point lies. The limit towards 11 is the one in which the $W$ and $Z$ become massless and a dimension opens. This note asks what that statement can mean.

It needs a path. On the interval of \cite{boundarynote}, the vector masses obey $\sin^2(M_Wl)=\tfrac12\sin^2 2\theta$, so $\theta\to0$ at fixed $l$ makes them vanish with the tower spacing fixed. The criterion we use is a hierarchy along a family indexed by $\varepsilon\to0$,
\begin{equation}\label{eq:hierarchy}
m_\varepsilon\ll E_\varepsilon\ll\Delta_\varepsilon\to0,
\end{equation}
with $m_\varepsilon$ the light scale and $\Delta_\varepsilon$ the first discarded level. Then a window at the moving scale $E_\varepsilon$ holds a fixed number of states, while the count at any fixed $E>0$ diverges. The two observations have different limits, and this reconciles a finite description near the weak scale with a tower at fixed energy.

The results of Sections~\ref{sec:interval} and \ref{sec:cylinder}, and the formulation of \eqref{eq:hierarchy}, follow an outside audit of this project (see the Provenance). The new content is Section~\ref{sec:window} and what rests on it: a window obtained from a local operator on a connected seven-dimensional parent, rather than prescribed.

\section{The interval}\label{sec:interval}

Let $(H_F,D_*,J_F,\Gamma_F)$ have the signs of KO-dimension six, with $r=\dim H_F$. For $0<\alpha<1$, let $D_F=\varepsilon D_*$ and $l_\varepsilon=l_*\varepsilon^{-\alpha}$, and let $D_\varepsilon$ be the derivative on the interval with the boundary conditions that realize $D_F$ exactly \cite{boundarynote}. Its signed spectrum is $(n\pi\pm\theta_k)/l$, $n\in\mathbb Z$, with $\theta_k=l\sigma_k$ and $\sigma_k$ the singular values of the block of $D_F$.

\begin{theorem}\label{thm:interval}
For small $\varepsilon$, an isometry $W_\varepsilon$ gives $W_\varepsilon^*D_\varepsilon W_\varepsilon=\varepsilon D_*$, with $J_F$ and $\Gamma_F$ exact, and
\begin{equation}
\begin{aligned}
\Delta_\varepsilon&\ge\frac{\pi}{l_*}\varepsilon^\alpha-\varepsilon\|D_*\|,\\
\|Q_\varepsilon(a)-\pi(a)\|&\le\tfrac23l_*^2\|D_*\|^2\varepsilon^{2(1-\alpha)}\|a\|.
\end{aligned}
\end{equation}
For $\alpha<\beta<1$ and $E_\varepsilon=E_*\varepsilon^\beta$, the number of levels with $|D_\varepsilon|\le E_\varepsilon$ is $r$. For fixed $E>0$ it is asymptotic to $(2rl_*E/\pi)\,\varepsilon^{-\alpha}$.
\end{theorem}

This is a direct consequence of the exact formulas (\supp{30}): $l_\varepsilon\|D_F\|\to0$, and $\varepsilon\ll\varepsilon^\beta\ll\varepsilon^\alpha$. The finite data remain inputs, the wall data vary by hand, and a fixed Majorana block does not scale as $\varepsilon D_*$, so an application to an electroweak Dirac block needs a reducing sector.

In the five-dimensional Yang-Mills toy of \cite{boundarynote}, with universal mass $m_\varepsilon=\varepsilon m_*$ and $\theta_\varepsilon=l_\varepsilon m_\varepsilon\to0$, the light poles are $M_W=l^{-1}\arcsin(\sin2\theta/\sqrt2)$ and $M_Z=l^{-1}\arcsin(2\sin2\theta/\sqrt5)$, so
\begin{equation}
M_W^2\simeq2m_\varepsilon^2,\qquad M_Z^2\simeq\tfrac{16}{5}m_\varepsilon^2,
\end{equation}
and the vector poles follow \eqref{eq:hierarchy}. The small-mass coefficients of the finite spectral action in the fully massive family are $3$ and $24/5$. The sum rule of \cite{boundarynote}, $M^2=\tfrac12(3-d_{\mathrm{eff}}(s))\,\|[D_F,T]\|^2/\|T\|^2$ with $s=l\Lambda$, connects the two pairs: $d_{\mathrm{eff}}=0$ gives the finite values, and $d_{\mathrm{eff}}=1$ the interval values. Along $\Lambda=E_\varepsilon$ we have $s\to0$, and at fixed $\Lambda$ we have $s\to\infty$. Varying the cutoff of the action is an additional prescription.

A family closer to the geometry below uses localization instead of boundary angles. With a bulk mass $\mu_\varepsilon=\mu_*\varepsilon^\alpha$ and $c_\varepsilon=\mu_\varepsilon l_\varepsilon=(1-\alpha)\log(1/\varepsilon)+c_0$, the light pair has $m=\mu\,\mathrm{sech}\,\kappa l$ with $\kappa l=c\tanh\kappa l$, so that
\begin{equation}
m_\varepsilon\sim2\mu_*e^{-c_0}\varepsilon,\qquad\Delta_\varepsilon\sim\mu_*\varepsilon^\alpha,\qquad\frac{1/\kappa}{l_\varepsilon}\to0.
\end{equation}
The widths and the length both grow, with a shrinking ratio. Writing the profile as $e^{-\kappa x}-re^{\kappa x}$, the opposite-wall coefficient is exactly $r=(m/(\mu+\kappa))^2$, quadratic in the light mass.

\section{A cylinder with a prescribed window}\label{sec:cylinder}

Let $A$ be the Dolbeault-Dirac operator of $\CP^3$ twisted by $\mathcal O(1)$, with kernel $\C^4$ and first nonzero level $\Delta_K$, and let $\mathcal D_r=E\partial_x+ZA/r$ on $I_l\times\CP^3$, with the spectral walls of \supp{24}. Taking $r_\varepsilon=r_*\varepsilon^{-\gamma}$ and the family above, the next level is at least $\min\{(\Delta_K/r_*)\varepsilon^\gamma,\,(\pi/l_*)\varepsilon^\alpha-\varepsilon\|D_*\|\}$, so a cutoff $E_*\varepsilon^\beta$ with $\max(\alpha,\gamma)<\beta<1$ isolates the window, and the count at fixed energy grows as $l_\varepsilon r_\varepsilon^6E^7$. For $\gamma=\alpha$ all seven directions expand homothetically. So \eqref{eq:hierarchy} is compatible with seven internal dimensions, for $4+7=11$. The window here is imposed by nonlocal walls.

\section{A window from a local operator}\label{sec:window}

Let $X$ be the disc bundle of $\Lambda^2_-T^*S^4$, with $G=Sp(2)$, $K=Sp(1)_a\times Sp(1)_b$, and $H=Sp(1)_a\times U(1)_b$. It has zero section $S^4$ and boundary $\CP^3$, and a smooth function has independent values at the two ends. Let $E=S^3(S)$, and let $J$ be the spin-$\frac32$ action on it. On $X$ take
\begin{equation}
T_M=D_X^E+i\Phi_M,\qquad\Phi_M=\epsilon Mf(r)\,n\cdot J,
\end{equation}
with $\hatD_M$ the doubled operator, $Z=-ic(dr)$, and the local wall condition $Zu=-u$ for $T_M$, $Zv=+v$ for $T_M^*$, which gives a self-adjoint elliptic operator (\supp{28}, Theorem 28.9). The sign $\epsilon$ is the one that puts the kernel of $T_M$ at the zero section. We assume the following: the metric is $dr^2+\pi^*g_{S^4}+\rho(r)^2g_{S^2}$, with $\rho=r$ and $f=r$ near $0$, and $\rho=1$, $f=f_0$ on a product collar, which is the Fubini-Study metric on $\CP^3$; the connection on $E$ is the pull-back near $S^4$; on the collar it is either (b) the connection preserving the four line summands $\mathcal O(3),\mathcal O(1),\mathcal O(-1),\mathcal O(-3)$, or (a) the pull-back.

\begin{theorem}\label{thm:window}
For $M$ large, $\hatD_M$ has exactly eight eigenvalues in $(-\delta,\delta)$, for every $\delta<\sqrt3$ in case (b), and for every $\delta<\min(2,\kappa_a)$ in case (a), with $\kappa_a\ge\sqrt5-\tfrac32$. The window is one copy of the $\mathbf 4$ of $Sp(2)$ in each $\hat\Gamma$ degree, at the zero section in the domain of $T_M$ and at the wall in the domain of $T_M^*$.
\end{theorem}

The proof (\supp{29}) is a partition into three regions, and each estimate holds on a fixed manifold, so uniformity over the representations of $Sp(2)$ needs no separate argument.

\emph{The zero section.} Sections on a tube are $K$-equivariant maps $G\to\Sigma_4\otimes L^2(\R^3;\C^2\otimes\C^4)$. Put $D_H=i\sum\gamma_ae_a$ and $\hat Q=\gamma_5\sigma\cdot\partial_y+\epsilon My\cdot J$. Then $T_M=T_0+R$, with $T_0=D_H+i\hat Q$ and $|R|\le C|y|$ the O'Neill term. The operator $\hat Q$ acts on the values and is $K$-equivariant, so it commutes with the derivatives $e_a$. With the Clifford signs, this gives exactly
\begin{equation}\label{eq:square}
T_0^*T_0=D_H^2+\hat Q^*\hat Q.
\end{equation}
The fibre operator has kernel $(1,\mathbf 3)\oplus(1,\mathbf 1)$ and a gap of order $M$ above it, and on the kernel $D_H$ is a twisted Dirac operator on $S^4$, with $D_H^2=\mathrm{Cas}_G(V)+\tfrac32-\mathrm{Cas}_K$ \cite{Parthasarathy1972}. This vanishes only on the $\mathbf 4$ and is at least $4$ otherwise. Since \eqref{eq:square} also bounds $M^2\||y|u\|^2$, the term $R$ is a small perturbation. The stronger localization theorem of \cite{Maridakis2021} has a hypothesis that fails for this mass, since $(n\cdot J)^2$ depends on $n$, and it is not needed here; nor is the block analysis of \cite{BL1991}.

\emph{The bulk.} For $w$ supported away from both ends, $\|T_Mw\|^2=\|Dw\|^2+\|\Phi_Mw\|^2+\langle w,i[D,\Phi_M]w\rangle\ge(M^2f^2/4-CM)\|w\|^2$, for every representation.

\emph{The wall.} On the collar in case (b), with $q=\epsilon n\cdot J$ and $m=Mf_0$,
\begin{equation}
\begin{aligned}
\|T_Mw\|^2={}&\|(\partial_x+mq)w_+\|^2+\|Aw_+\|^2\\
&+\|(\partial_x-mq)w_-\|^2+\|Aw_-\|^2,
\end{aligned}
\end{equation}
where $A$ is the sum of line-twisted Dirac operators. Only the channels $Z=-1,\ q>0$ for $T_M$ and $Z=+1,\ q>0$ for $T_M^*$ have a stable profile, and on these the lower bound comes from $A$. On the Fubini-Study metric $D_k^2=C_2+\tfrac32-\tfrac38k^2$ on the $SU(4)$-type with Casimir $C_2$; the normalization passes the Kirchberg limit $R/3$ for $k=0$ \cite{Kirchberg1986}. So $\mathcal O(3)$ has as kernel exactly one copy of the $\mathbf 4$, in $Z=+1$, and all other values have $D^2\ge3$. In case (a) the line summands mix through a bounded term with overlap $\sqrt3/2$ between the profiles, and a Schur complement of order $1/m$ reduces the wall to an elliptic operator on $\CP^3$, with fibre part $w(\rho)=2(\rho^2+2)^2/(9+2(\rho^2+2)^2)$ for its kernel, so $w=2/3$ at the Fubini-Study radius. At other radii the theorem holds under one kernel condition.

The transition regions lie in the bulk, where the mass controls $\|u\|$ with a factor $M^{-2}$, so the partition costs nothing at the scale of the gap. The limit values were fixed by these estimates before a direct check in Python (Table~\ref{tab:values}).

\begin{table}
\caption{\label{tab:values}Lowest singular values of $T_M$ on the $Sp(2)$-representation $V$, case (b), from a radial solve in Python at $M=32$, $l=3$, against the limits fixed by Theorem~\ref{thm:window}. The smallest value on the $\mathbf 4$ is a mesh floor at this resolution and is omitted; the wall values have a mesh error of about $0.005$.}
\begin{ruledtabular}
\begin{tabular}{lll}
$V$&$M=32$&limit\\
\hline
$\mathbf 4$&1.7323, 1.9991, 2.4497&$\sqrt3,\ 2,\ \sqrt6$\\
$\mathbf{16}$&2.2351, 2.2409, 2.4497&$\sqrt5,\ \sqrt5,\ \sqrt6$\\
$\mathbf{20}$&2.2409, 2.8275, 2.8286&$\sqrt5,\ \sqrt8,\ \sqrt8$\\
$\mathbf{40}$&3.4626, 3.4642, 3.4672&$\sqrt{12}$ (three)
\end{tabular}
\end{ruledtabular}
\end{table}

\begin{proposition}\label{prop:tunnel}
Let $T_Mu=\lambda v$, $T_M^*v=\lambda u$ be a unit singular pair of the window. For radial $\chi$ equal to $1$ near $S^4$ and $0$ near the wall,
\begin{equation}
\lambda\bigl(\langle u,\chi u\rangle-\langle v,\chi v\rangle\bigr)=\langle v,c(d\chi)u\rangle,
\end{equation}
and for every $\eta\in(0,1)$, $|\lambda_M|\le C_\eta\exp\bigl(-(1-\eta)\tfrac M2\int f\,dr\bigr)$, the integral running over the bulk.
\end{proposition}

The identity is an integration by parts, and Agmon estimates from each end give the bound, with rate set by the smallest eigenvalue $\frac12$ of $|n\cdot J|$. The summands of $E$ on $\CP^3$ have indices $4,0,0,-4$, of total zero, as equivariant cobordism invariance requires since $E$ extends over $X$ (\cite{Witten1983}, preprint pp.~43-44). The window has index $\mathbf 4-\mathbf 4=0$, so index theory allows $\lambda\neq0$, and the eight states come from the estimates. The $\mathbf 4$ is pseudoreal, so this says nothing yet about the complex representations of the Standard Model.

\begin{theorem}\label{thm:compression}
For $a\in C^1(X)$, the compression to the window satisfies
\begin{equation}
\|Q_M(a)-(Q_B(a|_{S^4})\oplus Q_Y(a|_{\CP^3}))\|\le C\|a\|_{C^1}M^{-1/2},
\end{equation}
with $Q_B(a)=\int a(x)(1-x\cdot\Gamma)\,d\mu$, and with $Q_Y$ the level-one Toeplitz map of the twistor space in case (b), and its $w(\rho)$-weighted form in case (a). Real functions on $X$ generate $M_2(\H)\oplus M_4(\C)$ under the limit map.
\end{theorem}

This is a statement about compression maps. Generated algebras do not converge when compression maps do, in line with the setting of \cite{CvS2021}.

\section{The joint limit on the disc bundle}\label{sec:joint}

Under $g\to R^2g$, with the potential divided by $R$, the operator is unitarily equivalent to $\hatD_M/R$. So the light mass and the gap both scale as $1/R$, with the compression maps unchanged. Let $G_0=\int f\,dr$ over the bulk, $0<\alpha<\beta<1$, and $\eta\in(0,1)$.

\begin{theorem}\label{thm:joint}
Take $M_\varepsilon=\frac{2(1-\alpha)}{(1-\eta)G_0}\log\frac1\varepsilon$, $R_\varepsilon=\varepsilon^{-\alpha}$, and $D_\varepsilon=\hatD_{M_\varepsilon}/R_\varepsilon$, in the setting of Theorem~\ref{thm:window}. For small $\varepsilon$:
\begin{enumerate}
\item $D_\varepsilon$ has exactly eight eigenvalues in $(-\delta\varepsilon^\alpha,\delta\varepsilon^\alpha)$, equal to $\pm m_\varepsilon$ with multiplicity four, and $m_\varepsilon\le C\varepsilon$;
\item the number of levels with $|D_\varepsilon|\le E_*\varepsilon^\beta$ is $8$;
\item the compression converges to the endpoint maps, with error $O((\log1/\varepsilon)^{-1/2})$;
\item at fixed $E>0$, the number of levels with $|D_\varepsilon|\le E$ is asymptotic to $C_XE^7\varepsilon^{-7\alpha}$.
\end{enumerate}
\end{theorem}

Item 4 follows from the Weyl law for $\hatD_0$, since $\hatD_M-\hatD_0$ is bounded by $cM_\varepsilon$, which is small against $ER_\varepsilon$. So \eqref{eq:hierarchy} holds for a window that comes from the operator, with growth of a seven-dimensional internal space at fixed energy. The bound on $m_\varepsilon$ is one-sided, since a lower bound on $\lambda$ is open. The fixed shape expands, so the pointed limit at fixed energy is locally flat at a generic point. Nothing here identifies it with the complete Bryant-Salamon manifold \cite{BS1989,KL2020}, whose metric lies outside our class, or with a solution of field equations.

\section{The quaternionic defect}

Proposition 28.16 of \supp{28} gives the test for the core block: with $B(r)=F^*F=\mathrm{diag}(b_aI_2,b_+,b_-)$, all real multipliers compress to operators commuting with the quaternionic structure $j$ if and only if $b_+=b_-$. Put $\eta_t=\int|b_+-b_-|\,dr/\int\mathrm{tr}\,B\,dr$, and let $P_H(A)=\frac12(A+jAj^{-1})$.

\begin{proposition}\label{prop:defect}
For real $f$ with $\|f\|_\infty\le1$, $\|Q(f)-P_HQ(f)\|\le2\eta_t$, and products of $k$ such compressions have defect at most $2k\eta_t$. Moreover, for every $\eta\in(0,1)$,
\begin{equation}
\eta_t\le C_\eta\exp\Bigl(-(1-\eta)M\int f\,dr\Bigr).
\end{equation}
\end{proposition}

The charge conjugation $C$ commutes with $T_M$ and exchanges the two wall conditions, so only the wall breaks the symmetry. On $4$-copies, $F\mapsto CFj^{-1}$ is a real structure, which fixes the model kernel. Then $\zeta=u-Cu$ obeys $T_M\zeta=\lambda(v-Cv)$ in the interior, with no component on the model kernel, and it is exponentially small at the core, sourced at the wall. A weighted Cauchy-Schwarz inequality gives twice the rate of Proposition~\ref{prop:tunnel}. So bounded expressions in the core block lie exponentially close to $M_2(\H)$. An unrestricted generated algebra may use coefficients of size $1/\eta_t$, and it is $M_4(\C)$ at finite mass when $b_+\ne b_-$ and the $\mathbf 5$ component is nonzero. The wall block generates $M_4(\C)$, which requires the loss of $C$ there.

\section{What is missing}

\emph{The vector sector.} On the family of Theorem~\ref{thm:joint} one needs the kinetic matrix and the mass Hessian of the vector fluctuations, then the momentum-dependent kernel, with the aim $M_W,M_Z\sim c\lambda/R$. The tunnelling eigenvalue alone does not establish this, so the words ``massless electroweak limit'' apply to the disc bundle only as a target.

\emph{The couplings.} For a bulk Yang-Mills field with constant profile, $g_4^{-2}=\mathrm{Vol}(X)/g_{11}^2$, so a plain dilation drives $g_4$ to zero, which must be separated from restoration at finite coupling. Gauge fields from isometries have Killing overlap integrals instead, and non-Killing fields behave differently again \cite{Baptista2026}.

\emph{The dynamics.} A potential $a\,r^{-p}+b\,\varepsilon^2r^q$ has its minimum at $r\propto\varepsilon^{-2/(p+q)}$, so it would select $\alpha=2/(p+q)$ if $p+q>2$. Nothing has derived such a potential, and adding it by hand derives nothing.

\emph{The finite structure.} The window gives the graded eight-dimensional space and the limit algebra. The connector of the 32-dimensional real bimodule, the quaternionic grading and the chirality of the physical spectrum remain inputs. The product of two copies of the parent has heat exponent $t^{-7}$, so the bimodule should come from a finite-rank correspondence, to keep seven dimensions.

\section{Status}

Theorem~\ref{thm:interval}, the vector poles and the cylinder are direct consequences of refereed lemmas, with proofs in \supp{30}. Theorems~\ref{thm:window}, \ref{thm:compression} and \ref{thm:joint} and Propositions~\ref{prop:tunnel}, \ref{prop:defect} are new here, for the stated class of metrics, with proofs in \supp{29} and \supp{30}. At this draft they are checked only by their author, and an independent referee is queued. Other metrics, and case (a) at other radii, stay open. Table~\ref{tab:values} is a Python computation, with the script in the repository.

\section*{Provenance}

The formulation of the joint limit and the interval and cylinder statements come from an outside audit of the repository, of unstated model, kept in its \texttt{reviews} folder. The proofs of Sections \ref{sec:window} and \ref{sec:joint} and of the defect bound are by Claude Opus 5.5, using results of GPT-6 Astra on the operator and its two ends, with referee reports by GPT-6.1 Terra on the earlier lemmas. Claude Sonnet 5.5 rechecked the audit in Python. Referee status is recorded in the repository ledger.

\bibliography{towerless_intervals}

\end{document}
